中国船舶当前不缺多头叙事，缺的是足以支持新增仓位的估值缺口。33.02元现价与33.08元独立基础价值几乎重合，多锚最终目标价33.07元仅对应约+0.15\%预期回报；因此我们的动作不是“追随周期买入”，而是\textbf{中性偏多（持有/等待验证）}。已有仓位可保留观察敞口，新增资金等待正式半年报同时证明利润、核心业务毛利率和经营现金流。\srcid{PR-01, CF-05}

\begin{keyinsight}[投委会一句话结论]
订单存量与2026年利润弹性已经得到正式披露支持，但当前价格也已为这些优势支付了相当信用。下一步的“所以呢”是：不以订单总额、重组标签或军工想象追价；只有半年报质量与2027年持续性证据使基础价值相对现价产生至少15\%上行时，才升级为新增仓位。
\end{keyinsight}

\section{E01｜一页决策仪表盘}

先看估值和动作，再看支撑动作的条件。基础公允区间是31.64--35.96元，反映2027E基础EPS 2.8767元对应11.0--12.5倍周期正常化PE；11.21元和48.34元是完整熊、牛经营情景，并不是基础区间的上下限。

\begin{exhibitbox}[核心估值与动作]
\small
\begin{tabularx}{\exhibitboxwidth}{L{2.7cm}R{2.2cm}R{2.4cm}R{2.2cm}X}
\toprule
\textbf{标的} & \textbf{现价} & \textbf{最终目标/区间} & \textbf{预期回报} & \textbf{评级与动作} \\
\midrule
中国船舶\newline 600150.SH & 33.02元 & 33.07元\newline 31.64--35.96元 & +0.15\% & 中性偏多；已有仓位持有，新增资金等待正式半年报验证 \\
\bottomrule
\end{tabularx}
\vspace{0.35em}
\begin{tabularx}{\exhibitboxwidth}{L{2.7cm}R{2.2cm}R{2.4cm}R{2.2cm}X}
\toprule
\textbf{估值母体} & \textbf{熊市价值} & \textbf{基础价值} & \textbf{牛市价值} & \textbf{方法/失效条件} \\
\midrule
仅600150合并归母 & 11.21元 & 33.08元 & 48.34元 & 2027E周期正常化PE；利润、毛利和现金任一显著失速即下修 \\
\bottomrule
\end{tabularx}
\sourcenote{现价PR-01；股本CF-04；估值为冻结R1模型。}
\end{exhibitbox}

该表的投资含义很直接：基础情景没有显著安全边际，熊市情景却显示非对称下行。牛市价值并非“目标价”，它要求交付、价格/结构、毛利和现金转化同时处于上端。组合行为应以事件验证为主，而不是按单一高景气标签配置。

\section{下一验证点：正式半年报而非预告年化}

2026H1归母净利润预告为92亿--110亿元，扣非为90亿--108亿元，性质是未经审计业绩预告；收入、毛利率、经营现金流与营运资金尚未披露，不能将区间中点乘二得到全年预测。2026Q1归母48.32亿元后，若H1达到中点101亿元，则Q2需贡献52.68亿元；若扣非达到中点99亿元，则Q2需贡献51.33亿元。\srcid{CF-04, CF-05}

\begin{exhibitbox}[下一季度盈利与质量门槛]
\small
\begin{tabularx}{\exhibitboxwidth}{L{3.0cm}C{2.4cm}C{2.4cm}X}
\toprule
\textbf{指标} & \textbf{正向门槛} & \textbf{降级门槛} & \textbf{动作含义} \\
\midrule
H1归母/扣非 & 至少101/99亿元 & 低于92/90亿元 & 高于中点且质量过关才进入升级评估 \\
核心业务毛利率 & 不低于11.72\% & 低于10.5\% & 调整2026--2028毛利假设与估值倍数 \\
经营现金流 & 为正，且营运资金无异常扩张 & 为负且存货/合同资产快于收入 & 现金背离时维持或降级，不因利润高增追价 \\
年度交付 & 接近168艘/1,650万DWT计划 & 重大延期或计划明显失速 & 下修收入确认与持续期 \\
全年OCF/归母 & 至少0.8倍 & 低于0.8倍且无回收路径 & 取消11.5倍基础倍数信用 \\
\bottomrule
\end{tabularx}
\sourcenote{CF-03、CF-04、CF-05、CF-17；阈值为冻结House与R1模型。}
\end{exhibitbox}

这组门槛的“所以呢”是：利润绝对额只是第一层，毛利率说明高价订单是否真正赚钱，现金流说明利润是否可回收。三者缺一时，都不足以把“持有等待验证”升级为“新增仓位”。

\section{E16｜概率情景与组合纪律}

我们不把熊、基、牛价值简单平均为目标价，因为三者代表不同的经营状态和证据门槛。基础状态下，估值回报几乎为零；投资收益的主要来源必须是正式结果推动盈利预测上修，而不是依赖倍数扩张。

\begin{exhibitbox}[情景、条件与监测等价行为]
\small
\begin{tabularx}{\exhibitboxwidth}{L{2.0cm}R{2.0cm}L{4.8cm}X}
\toprule
\textbf{情景} & \textbf{价值} & \textbf{成立条件} & \textbf{监测等价行为} \\
\midrule
熊市 & 11.21元 & 交付/确认走弱，核心毛利受成本、返工或低价旧单侵蚀，现金转化偏弱 & 失效条件触发即降低观察敞口并重跑模型 \\
基础 & 33.08元 & 高价订单按计划交付，2027E EPS 2.8767元，现金逐步改善 & 持有；区间中部附近不追价 \\
牛市 & 48.34元 & 交付超计划，价格/结构与效率同时改善，OCF同步处于上端 & 仅在正式数据连续验证后逐步提高研究权重 \\
\bottomrule
\end{tabularx}
\sourcenote{冻结增长模型与估值模型；情景值按2027E EPS乘对应PE。}
\end{exhibitbox}

风险标签“高”指一旦发生足以使基础价值、盈利分母或估值倍数发生重大下修，并不等于事件发生概率高。监测等价行为也不是交易指令：核心审查名单仅包含600150；海外同业、中国动力与沪东中华仅是竞争或边界观察节点，不分享600150的目标价。

\section{边界、最强反方与主张失效}

报告全程只对一个估值母体定价：600150.SH合并口径。江南造船、大连造船、外高桥造船、广船国际、武昌造船、中船澄西和北海造船都是合并组成，不进行子公司SOTP；中国动力是参股观察节点；沪东中华位于上市体外。未披露军品、集团未分配订单和潜在资产注入全部零计价。中远海运约500亿元、87艘项目已在上市公司聚合新签/在手订单中反映，且含部分意向订单，不能再次加入收入、EPS或目标价。\srcid{CF-03, CF-12}

最强反方是：等待正式半年报可能错过高确定性的盈利上修。H1预告、在手/收入约3.08倍以及2026交付计划，确实可能使2027年盈利继续上修。我们接受这一反方，但它成立必须同时看到H1靠近上端、核心毛利不回落、经营现金流不背离以及交付按计划；只满足利润一项，不足以证明高利润持续期。

\begin{riskbox}[什么会证明本报告错了]
若正式H1归母/扣非达到或超过110/108亿元，核心毛利显著高于11.72\%，经营现金流同步改善，并出现可审计的2027--2028高价订单排程，则本报告可能过于谨慎，应快速重跑盈利与估值。反之，H1低于92/90亿元、核心毛利低于10.5\%、重大延期/撤单/客户融资问题或全年OCF/归母低于0.8倍，将证明市场多头主张过度乐观。
\end{riskbox}

最终动作由证据而不是标签改变：正式披露出现时，优先重跑合并收入、Growth毛利、少数股东归属与现金转化；在此之前，订单存量、军工能力、集团排名和弱来源目标价都不能替代估值安全边际。
